\documentclass[11pt]{article}
\usepackage[a4paper,left=26mm,right=34mm,top=22mm,bottom=22mm]{geometry}
\usepackage{fontspec,tikz}
\usepackage[unicode,hidelinks]{hyperref}
\usepackage{polyglossia}
\setmainlanguage{english}
\setotherlanguages{arabic}
\setmainfont{Linux Libertine O}[Ligatures=TeX]
\newfontfamily\arabicfont{Amiri}[Script=Arabic]
\newfontfamily\ArBody{Amiri}[Script=Arabic,Scale=1.08]
\newfontfamily\ArApp{Amiri}[Script=Arabic,Scale=0.92]
\newfontfamily\ArC{Amiri}[Script=Arabic]
\hypersetup{pdftitle={al-Battani, Kitab fi dala'il al-qiranat wa-l-kusufat (astrological history), ed. G. Saliba (Suhayl 9). S15},pdfauthor={al-Battani; ed. George Saliba}}
\setlength{\parindent}{0pt}\setlength{\parskip}{0pt}
\tracinglostchars=2
\newcommand{\PageHead}[3]{\par\bigskip\penalty-300\pdfbookmark[1]{p. #2 (PDF #1)}{pg-#1}\hypertarget{AB02-PDF#3}{}%
 {\large\bfseries p.\,#2}\enspace{\footnotesize (PDF #1)}\par\smallskip}
% one printed line: right to left; in the right margin its number and, where printed, the folio of the manuscript
\newcommand{\EdL}[4]{\par\noindent\hypertarget{#2}{}\hbox to\linewidth{\parbox[t]{\linewidth}{\setRTL\ArBody\raggedleft #4\par}\rlap{\hspace{3mm}\footnotesize #1\enspace #3}}\par}
\newcommand{\EdLc}[4]{\par\noindent\hypertarget{#2}{}\hbox to\linewidth{\parbox[t]{\linewidth}{\setRTL\ArBody\centering #4\par}\rlap{\hspace{3mm}\footnotesize #1\enspace #3}}\par}
\newcommand{\EdLh}[4]{\EdLc{#1}{#2}{#3}{\bfseries #4}}
\newcommand{\EdLp}[4]{\EdL{#1}{#2}{#3}{\hspace*{6mm}#4}}
\newcommand{\AppL}[2]{\par\noindent\hypertarget{#1}{}\parbox[t]{\linewidth}{\setRTL\ArApp\raggedleft #2\par}\par}
\begin{document}

\thispagestyle{empty}
\begin{center}\Large AL-BATTĀNĪ\par\vspace{4mm}\LARGE The book on the indications of the conjunctions and the eclipses\par\vspace{2mm}\large(his astrological history)\par\vspace{4mm}\large\RL{\ArBody كتاب في دلائل القرانات والكسوفات}\par\vspace{5mm}\normalsize Ankara, Dil ve Tarih-Coğrafya Fakültesi, collection Saib, 199/2, ff.\,27v–42v\par\vspace{4mm}The Arabic text as edited by George Saliba, in E.\,S. Kennedy, B. van Dalen, G. Saliba and J. Samsó,\\«Al-Battānī’s Astrological History of the Prophet and the Early Caliphate», \textit{Suhayl} 9 (2009–2010), pp.\,13–148: the Arabic text on pp.\,18–78\par\vspace{4mm}\large S15 --- the Arabic text, line by line\end{center}\vspace{6mm}
\noindent The edition of the Arabic text is printed in the article as images of George Saliba’s typeset edition, on the even pages 18–78. This file transcribes those pages: every printed line of the text is one line here, numbered in the right margin as the editors number it, with the folio of the manuscript where the edition prints it; the horoscope charts are transcribed cell by cell and drawn again; the editors’ apparatus follows each page. The justification of the print (the long joins of the letters) is not transcribed; the tanwīn of the accusative is written on the letter before the final alif (on the lām-alif in \RL{\ArBody لاً}) wherever the print sets it; lines printed with a paragraph indent are indented here too. The English translation and the commentary of the article are not reproduced. The article is licensed CC BY-NC-ND 3.0 ES, and this transcription is distributed under the same licence, with the attribution above. The transcription is AI-integrated work (Claude); no person has reviewed it.\par\clearpage
\PageHead{6}{18}{0006}
\EdLc{1}{AB02-PDF0006-L01}{\RL{\ArBody ٢٧ظ}}{/بسم اللّه الرّحمن الرّحيم}
\EdL{2}{AB02-PDF0006-L02}{}{قال محمد بن جابر البتّاني :}
\EdL{3}{AB02-PDF0006-L03}{}{لما كان النظر في أمور الملل وأوقات الدول وأحوال الملوك ومبادئ ذلك وأوقات تغييراته من أجلّ}
\EdL{4}{AB02-PDF0006-L04}{}{ما تقدمت به المعرفة من صناعة النجوم ، وكُنّا قد وقفنا على ما تكلم به في ذلك كثيرٌ من أهلها في زماننا ،}
\EdL{5}{AB02-PDF0006-L05}{}{فوجدناهم قد راموا علمها والتمسوا معرفتها بتحاويل السنين التي جعلوا مبادئها دخول الشمس برج الحمل}
\EdL{6}{AB02-PDF0006-L06}{}{في أزمان القرانات بطوالع تلك الأوقات ، وأشياء كثيرة من هذا النوع ليس فيها حجّة توجب}
\EdL{7}{AB02-PDF0006-L07}{}{استعمالها ولا أصول يقاس عليها ، وأغفلوا أمر الكسوفات}
\EdL{8}{AB02-PDF0006-L08}{}{الواقعة مع القرانات في مواضعها من فلك البروج واشتراكات الكواكب في الأشكال ،}
\EdL{9}{AB02-PDF0006-L09}{}{إذ كان ذلك من أدلّ الأشياء على التغييرات وأوّل الانقلابات ، وكان بطلميوس ، مع تقدُّمه}
\EdL{10}{AB02-PDF0006-L10}{}{في هذه الصناعة ، لم يعوّل علي شيء غيرها في معرفة الأحداث الكائنة في العالم ،}
\EdL{11}{AB02-PDF0006-L11}{}{رأينا أن نعتقد مذهبه ونتبع ما رسمه وأشار اليه من ذلك في كتابه المعروف بالأربع مقالات}
\EdL{12}{AB02-PDF0006-L12}{}{الذي ذكر فيه القضاء علي الحوادث من النجوم . فنقدم قبل وصفنا قولاً وجيزًا . وهو أنّ أثقل الكواكب}
\EdL{13}{AB02-PDF0006-L13}{}{المتحيرة وأبطأها مقارنة وأعلاها أفلاكًا وأعظمها أجرامًا كوكب زحل وكوكب}
\EdL{14}{AB02-PDF0006-L14}{}{المشتري . ولذلك خُصّا بالقران دون غيرهما من الكواكب الباقية . ويقع اقترانهما في زهاء}
\EdL{15}{AB02-PDF0006-L15}{}{كل عشرين سنة مرّة واحدة . فيقترنان في كل مثلّثة من مثلّثات البروج اثنتي عشرة}
\EdL{16}{AB02-PDF0006-L16}{}{مرّة ، في كل برج أربع مرّات . ثم تنتقل الى المثلّثة التى تتلوها على ذلك الرسم حتّى يتم اقترانهما}
\EdL{17}{AB02-PDF0006-L17}{}{في مثلّثات فلك البروج الأربع في مقدار تسع ماية وثلاث وخمسين سنة شمسيّة بالتقريب .}
\par\smallskip\noindent\rule{40mm}{0.4pt}\par
\AppL{AB02-PDF0006-A01}{٥- معرفتها : معرفته ، الشمس: بدأ الناسخ بكتابة كلمة «أوّل» بعد هذه الكلمة فكتب «أو» ثم عاد وشطبها | ٩- فلك: كرر الناسخ}
\AppL{AB02-PDF0006-A02}{سهوًا هذه الكلمة فكتب «الفلك فلك» | ٩- أدلّ: أجل | ٩ - وأوّل : كتب الناسخ «واوجها» وصحّحت على الهامش | ١٢ - الذى :}
\AppL{AB02-PDF0006-A03}{التي ، قبل: قبل ما | ١٤ - ولذلك: ولذلك ما | ١٦ - تنتقل: المقصود هي القرانات}
\PageHead{8}{20}{0008}
\EdL{1}{AB02-PDF0008-L01}{\RL{\ArBody ٢٨و}}{/ ثم يعود القران إلى البرج الخامس من البرج الذي اقترنا فيه أوّل مرّة ثم إلى البرج التاسع منه بعد}
\EdL{2}{AB02-PDF0008-L02}{}{ألف وتسع مائة وست وستين شمسيّة حتّى تكون عودته إلى البرج الأوّل بعد ألفي وثمان مائة}
\EdL{3}{AB02-PDF0008-L03}{}{سنة وتسع وخمسين سنة. وذلك بعد مائة وأربعة وأربعين قرانًا . فإذا ضرب كلٌّ واحد من هذين العددين}
\EdL{4}{AB02-PDF0008-L04}{}{في عدد البروج الإثني عشر كان الذي يجتمع من عدد القرانات}
\EdL{5}{AB02-PDF0008-L05}{}{ألف وسبع مائة وثمانية وعشرون قرانًا . ومن عدد السنين قريبًا من أربعة وثلاثين ألفًا}
\EdL{6}{AB02-PDF0008-L06}{}{وثلاثمائة سنة . هذا شبيه بالعدد الذي ذكرت أنّه مقدار دورة واحدة من أدوار النجوم . ولكوكب المريخ}
\EdL{7}{AB02-PDF0008-L07}{}{أيضًا مع هذين الكوكبين شهادة واشتراك قويّ لمكانه فوق الشمس ولأنّه أعظم جرمًا من الكوكبين الباقيين}
\EdL{8}{AB02-PDF0008-L08}{}{الذين تحتها ، أعني الزهرة وعطارد . وقوّتهما}
\EdL{9}{AB02-PDF0008-L09}{}{في ذلك دون قوّة المريخ . فكلما اقترن زحل والمشتري في مثلّثة ظهرت قوّتهما وعزّ أهل البلدان}
\EdL{10}{AB02-PDF0008-L10}{}{والأمصار والأمم المنسوبة إليهما وقهروا من يليهم ، سيّما ما كان من هذه البلدان منسوبًا إلى}
\EdL{11}{AB02-PDF0008-L11}{}{البرج الذي يكون فيه القران والكسوف جميعًا. وينالهم المضرّة والنكبات عند تردّد زحل والمريخ في برجهم}
\EdL{12}{AB02-PDF0008-L12}{}{أو في البرج المقابل له أو على تربيعه. ويكون ذلك بحسب الترتيب}
\EdL{13}{AB02-PDF0008-L13}{}{في القوّة والضعف. وكذلك ينالهم الخير والسعادة في الأمور عند تردد المشتري والزهرة}
\EdL{14}{AB02-PDF0008-L14}{}{في هذه الأوتاد وخاصّة المشتري لطول مقامه في الموضع. والقول فيما ذكرنا}
\EdL{15}{AB02-PDF0008-L15}{}{ينقسم إلى ثلاثة أقسام. فالقسم الأوّل أُمور الملل وهو أعلاها مرتبة وأطولها مدّة وأثقلها وزنًا.}
\EdL{16}{AB02-PDF0008-L16}{}{والقسم الثاني أُمور الدول الذي لها انتقالات في البيوتات من الملّة وهو أوسط الثلاثة في القوّة والثقل.}
\EdL{17}{AB02-PDF0008-L17}{}{والقسم الثالث أُمور الملوك وهو دون القسمين الأولين في القوّة والمدّة وأكثرها اختلافًا.}
\par\smallskip\noindent\rule{40mm}{0.4pt}\par
\AppL{AB02-PDF0008-A01}{١- البرج\textsuperscript{٢} : البروج | ٢ - وست : وستة | ٢- ٣ - ألفي ... سنة: كتب الرقم «٢٨٥٩» تحت السطر | ٣- مائة ... وأربعين: كتب}
\AppL{AB02-PDF0008-A02}{الرقم «١٤٤» تحت السطر | ٥ - ألف ... وعشرون: كتب الرقم «١٧٢٨» تحت السطر}
\PageHead{10}{22}{0010}
\EdL{1}{AB02-PDF0010-L01}{\RL{\ArBody ٢٨ظ}}{/وقد تكون الملّة بالدولة والدولة بالملّة وقد تضعف إحداهما أحيانًا وتقوى. فأمّا الدول}
\EdL{2}{AB02-PDF0010-L02}{}{التي تكون بغير ملل فإنّها سريعة الزوال. وكذلك الملل التي تكون بغير دول}
\EdL{3}{AB02-PDF0010-L03}{}{فإنّها تفسد وتضمحلّ سريعًا. ولزحل والمشتري والمرّيخ فضل على غيرها من الكواكب}
\EdL{4}{AB02-PDF0010-L04}{}{إذا كانت مقيمة عند الكسوف وكانت الدلالة مأخوذة منها باستيلائها على}
\EdL{5}{AB02-PDF0010-L05}{}{التدبير. ويكون استيلاؤها على التدبير بحسب ما لها من الحصص في المواضع من أنّها أرباب}
\EdL{6}{AB02-PDF0010-L06}{}{فلك البروج والشرف والمثلّثات والحدود والظهور القريب أو الاشتراك في الشكل،}
\EdL{7}{AB02-PDF0010-L07}{}{أعني في برج الكسوف وبرج الوتد الذي يجاريه أحد المضيئين قبل كسوفه . وإن اجتمعت}
\EdL{8}{AB02-PDF0010-L08}{}{هذه الحصص الخمس لكوكب واحدٍ بعينه كان هو صاحب التدبير وحده .}
\EdLp{9}{AB02-PDF0010-L09}{}{وإن كانت لكوكبين أو أكثر من ذلك ، أُخذ الكوكبان معًا اللذان لهما مشاكلات}
\EdL{10}{AB02-PDF0010-L10}{}{أكثر الى كلّ واحد من الموضعين أو ساير الكواكب التي تكون لها الحصص فإن}
\EdL{11}{AB02-PDF0010-L11}{}{كان ما لكل واحدٍ منها مساويًا للآخر في الوزن اختير وقُدَّم في التدبير ما كان أولى أن}
\EdL{12}{AB02-PDF0010-L12}{}{يُقال إنّه في الوتد . ويُختار ويُقدَّم أوّلاً صاحب تدبير البرج الذي فيه الكسوف . ويُؤخذ أيضًا}
\EdL{13}{AB02-PDF0010-L13}{}{مع ذلك من الكواكب الثابتة ما يطلع أو يتوسّط السماء مع موضع الكسوف ، أيّ}
\EdL{14}{AB02-PDF0010-L14}{}{هذين الوتدين اتفق مصير أحد المضيئين إليه وقت الكسوف ،}
\EdL{15}{AB02-PDF0010-L15}{}{وأيضًا الكولب التي تكون في وقت الكسوفات في ذلك الوتد . فيكون الاستدلال}
\EdL{16}{AB02-PDF0010-L16}{}{على خواصّ الأشياء من خواصّ البروج التي يكون فيها الكسوف ، والتي تكون فيها الكواكب}
\EdL{17}{AB02-PDF0010-L17}{}{المتحيّرة المتولّية للتدبير. ومن صور هذين الموضعين ومشاكلة الكواكب الثابتة}
\par\smallskip\noindent\rule{40mm}{0.4pt}\par
\AppL{AB02-PDF0010-A01}{٣ - ولزحل: كتب على هامش هذا السطر ، وبخط مختلف ، العبارة التالية «مدبّر كسوف وخسوف» | ٩ - اللذان: الذان |}
\AppL{AB02-PDF0010-A02}{١٤- الكسوف: هناك بياض في المخطوطة بعد هذه الكلمة ملئ بالعبارة التالية «مدبّر كسوف وخسوف من الثوابت» | ١٧-}
\AppL{AB02-PDF0010-A03}{المتولّية: المولّية}
\PageHead{12}{24}{0012}
\EdL{1}{AB02-PDF0012-L01}{\RL{\ArBody ٢٩و}}{/لتلك المتحيرة المستولية على التدبير في المزاج والقوة ، سيّما الكواكب التي هي أعظم}
\EdL{2}{AB02-PDF0012-L02}{}{ضياءً ، أعني التي تكون في القدر الأوّل ثم في القدر الثاني من الأقدار الستة التي رُسمت للكواكب}
\EdL{3}{AB02-PDF0012-L03}{}{الثابتة في عظم الأجرام ، وذلك أنّها إذا شاكلت الكواكب المتحيرة في المزاج والقوة}
\EdL{4}{AB02-PDF0012-L04}{}{قوّت دلائلها وصيّرتها محضة خالصة ، وكما أنّها إذا خالفتها في المزاج والقوة نقصت}
\EdL{5}{AB02-PDF0012-L05}{}{من تلك الدلائل ولم تكن تامّة في القوة ولا محضة ولا خالصة. وكذلك تفعل الكواكب}
\EdL{6}{AB02-PDF0012-L06}{}{المتحيرة إذا شاكلت الكواكب المتولّية للتدبير منها في الخير، أو كانت على خلاف [ ذلك ] .}
\EdLh{7}{AB02-PDF0012-L07}{}{دلائل الكواكب في البروج ونواحي الفلك}
\EdLp{8}{AB02-PDF0012-L08}{}{فإن كانت في أحد برجي الاعتدالين دلّت على ما يحدث في غلاّت الهياكل ومواضع العبادة}
\EdL{9}{AB02-PDF0012-L09}{}{والمساجد وما أشبهها. فإن كانت في أحد برجي الانقلابين دلّت على تغيّر السنن والعبادات ،}
\EdL{10}{AB02-PDF0012-L10}{}{نريد الأديان والملل. وإن كانت في ذوات الأجساد دلّت على ما يحدث في الناس وعلى الملوك .}
\EdL{11}{AB02-PDF0012-L11}{}{وإن كانت في البروج الثابتة دلّت على ما يحدث في الأساسات والأبنية، نريد الأُصول والشرايع.}
\EdL{12}{AB02-PDF0012-L12}{}{وكذلك أيضًا إذا كانت في المواضع التي وضعها في المشرق خاصّة ، نريد الطالع، دلّت على أنّ الحادث}
\EdL{13}{AB02-PDF0012-L13}{}{يكون في مَن سنّه سن الأحداث وفي الأساسات . وإن كانت في}
\EdL{14}{AB02-PDF0012-L14}{}{وسط السماء دلّت على الحادث أن يكون [في مَن ] سنه السنّ الوسطى وفي الهياكل والملوك .}
\EdL{15}{AB02-PDF0012-L15}{}{وفي ناحية المغرب دلّت على سنّ المشايخ النافذي الأعمار وعلى تغييرات السير أيضًا .}
\EdLp{16}{AB02-PDF0012-L16}{}{ونعلم كثرة الحادث وقلّته من قياس الكواكب المدبّرة إلى موضع الكسوف. وذلك}
\EdL{17}{AB02-PDF0012-L17}{}{أنّها متى كانت في أشكال منسوبات إلى العشيّات عند كسوفات القمر أو أشكال}
\par\smallskip\noindent\rule{40mm}{0.4pt}\par
\AppL{AB02-PDF0012-A01}{٦- ذلك: نقترح زيادتها لاستقامة المعنى | ٨- يحدث: يحث | ١٤-في مَن : نقترح إضافتها كما وردت سابقًا لاستقامة النصّ}
\PageHead{14}{26}{0014}
\EdL{1}{AB02-PDF0014-L01}{\RL{\ArBody ٢٩ظ}}{/ منسوبة إلى الغدوات عند كسوفات الشمس ، دلّت على كثرة الحادث وعلى أنّه يكون}
\EdL{2}{AB02-PDF0014-L02}{}{في أكثر الجنس . وإذا كانت مقابلة للشمس على القطر كان الحادث وسطًا}
\EdL{3}{AB02-PDF0014-L03}{}{وفي نصف ذلك الجنس . وإذا كانت في الأشكال المنسوبة إلى الغدوات عند كسوفات}
\EdL{4}{AB02-PDF0014-L04}{}{القمر، أو في الأشكال المنسوبة إلى العشيّات عند كسوفات الشمس كان الحادث}
\EdL{5}{AB02-PDF0014-L05}{}{قليلاً في أقلّ الجنس . ويستدلّ على سرعة الحادث وإبطائه من تلك الكواكب أيضًا .}
\EdL{6}{AB02-PDF0014-L06}{}{وذلك أنّها متى كانت في الأمرين جميعًا، أعني كسوفات الشمس والقمر معًا، منسوبة إلى}
\EdL{7}{AB02-PDF0014-L07}{}{الغدوات كان الحادث عنها سريعًا. ومتى كانت منسوبة إلى العشيّات كان}
\EdL{8}{AB02-PDF0014-L08}{}{عنها بطيئًا. وكذلك إذا كانت في ناحية المشرق من الأفق دلّت على السرعة.}
\EdL{9}{AB02-PDF0014-L09}{}{وإذا كانت ممّا يلي وسط السماء دلّت على التوسط في السرعة والإبطاء .}
\EdL{10}{AB02-PDF0014-L10}{}{وإذا كانت في ناحية المغرب دلّت على الإبطاء والتأخير. ومتى كان الكسوف}
\EdL{11}{AB02-PDF0014-L11}{}{ظاهر الحسّ متصاعدًا إلى وسط السماء كان الحادث قويًّا متزايدًا. فإذا كان منحدرًا إلى}
\EdL{12}{AB02-PDF0014-L12}{}{الأُفق صايرًا إلى المغرب والخفاء عن الحسّ كان دون ذلك .}
\EdLp{13}{AB02-PDF0014-L13}{}{والزهرة وعطارد إذا كانا منسوبين إلى طلوع العشيّات دلاّ على السرعة .}
\EdL{14}{AB02-PDF0014-L14}{}{وطلوعهما المنسوب [إلى] الغدوات يدلّ على دون ذلك . وغروبهما المنسوب إلى الغدوات}
\EdL{15}{AB02-PDF0014-L15}{}{والعشيّات جميعًا يدلّ على الإبطاء. وكثرة دلائلهما في كثرة الحادث وقلّته.}
\EdLp{16}{AB02-PDF0014-L16}{}{وأما دلالة كلّ واحدٍ من الكواكب الخاصّة له، فإنّ زحل إذا ولي التدبير وحده}
\EdL{17}{AB02-PDF0014-L17}{}{كان سببًا للخوف والضيق والهرب والموت وخاصّة لمن قد طعن في السن .}
\par\smallskip\noindent\rule{40mm}{0.4pt}\par
\AppL{AB02-PDF0014-A01}{١٢- الأُفق: من الأُفق}
\PageHead{16}{28}{0016}
\EdL{1}{AB02-PDF0016-L01}{\RL{\ArBody ٣٠و}}{/ وأما المشتري فإنّه يكون سببًا للسعادة والسلامة والعمارة والعدل ويزيد الخيرات}
\EdL{2}{AB02-PDF0016-L02}{}{وعطايا الملوك . وأما المرّيخ فإنّه يكون سببًا للحروب والسبي وتوثب}
\EdL{3}{AB02-PDF0016-L03}{}{جملة الناس وغضب الرؤساء ، سيما بين القوم وأهل جنسهم ، وما قد يصاب به بعض الناس}
\EdL{4}{AB02-PDF0016-L04}{}{بالغضب والهوان والتعدّي والخروج عن الشريعة والحريق والسلب والقتل .}
\EdL{5}{AB02-PDF0016-L05}{}{وأمّا الزهرة فإنّها إذا وليت التدبير وحدها كان ما يحدث عنها شبيهًا بما يحدث عن}
\EdL{6}{AB02-PDF0016-L06}{}{المشتري من رفع المراتب والكرامات وتحريك أمور الجماع وحُسن الأخلاق ومحبّة النظافة والطيب وهي}
\EdL{7}{AB02-PDF0016-L07}{}{دالّة على الوحي والأمور الإلهية . وأمّا عطارد فإنّه إذا مازج بعض الكواكب شاكل طبيعته، غير أنّه}
\EdL{8}{AB02-PDF0016-L08}{}{يحرك الأشياء أكثر من ساير الكواكب ، ويُعين على القوّة والجدّ في الأمور ولطف الحيلة وقطع الطرق}
\EdL{9}{AB02-PDF0016-L09}{}{والغارات وأمور الكهانة والعبادة وما يعرض في غلاّت الهياكل وتغيير السنن .}
\EdLh{10}{AB02-PDF0016-L10}{}{[ ألوان الكسوف ]}
\EdLp{11}{AB02-PDF0016-L11}{}{وأيضًا إذا كانت ألوان الكسوف سوداء أو مايلة إلى الخضرة دلّت على الأشياء التي من طبيعة زحل.}
\EdL{12}{AB02-PDF0016-L12}{}{وإذا كانت بيضاء دلّت على الأشياء التي من طبيعة المشتري . فإن كانت مايلة إلى الحمرة دلّت على}
\EdL{13}{AB02-PDF0016-L13}{}{الأشياء التي من طبيعة المرّيخ . وإذا كانت صهباء ذهبيّة دلّت على الأشياء التي من طبيعة الزهرة .}
\EdL{14}{AB02-PDF0016-L14}{}{وإذا كانت مختلفة الألوان دلّت على الأشياء التي من طبيعة عطارد .}
\EdLp{15}{AB02-PDF0016-L15}{}{وإذا كان اللون في جميع جرم أحد المضيئين أو فيما حوله كان ذلك في جميع أجزاء البلدان .}
\EdL{16}{AB02-PDF0016-L16}{}{وإذا كان في جزء من النيّرين أو فيما حولهما كان في الناحية التي تلي ذلك الجزء فقط .}
\EdLp{17}{AB02-PDF0016-L17}{}{وكذلك أيضًا يستدلّ بذوات الأذناب التي تظهر}
\par\smallskip\noindent\rule{40mm}{0.4pt}\par
\AppL{AB02-PDF0016-A01}{٦- من: مع | ٩- السنن: يلي هذه الكلمة العبارة التالية «الوان كسوف» ، أُدخلت في النص لاحقًا، أثبتناها هنا موافقة مع نص}
\AppL{AB02-PDF0016-A02}{ابن أماجور المثبت في الملحق .}
\PageHead{18}{30}{0018}
\EdL{1}{AB02-PDF0018-L01}{\RL{\ArBody ٣٠ظ}}{/وساير ما يحدث من [الأ]شياء التي تسمّى [ العصيّ] والبوقات وما أشبهها. وتُقام طبائعها مقام ما قيل}
\EdL{2}{AB02-PDF0018-L02}{}{إنّ طبيعة المرّيخ وعطارد تحدث وتظهر فيهما من الحروب والاحتراق والزلازل . ويُستدل على}
\EdL{3}{AB02-PDF0018-L03}{}{المواضع التي تحدث فيها الأحداث بمواضعها من فلك البروج وبالمواضع التي يسامتها}
\EdL{4}{AB02-PDF0018-L04}{}{ذنب الكوكب . ويستدل على ما يحدث من صورة ما يظهر، وفي زمان مدّته على مدّة}
\EdL{5}{AB02-PDF0018-L05}{}{الأحداث ومن قياسه على الوقت في كون الحادث . وذلك أنّها في أكثر الأمر إذا كان}
\EdL{6}{AB02-PDF0018-L06}{}{ظهورها بالعشيّات كان الحادث عنها بطيئًا . وإذا ظهرت بالغدوات كان الحادث}
\EdL{7}{AB02-PDF0018-L07}{}{سريعًا. وممّا يُستدل به على قوّة الحادث وضعفه أن تكون الكسوفات مع القرانات}
\EdL{8}{AB02-PDF0018-L08}{}{في مواضعها من فلك البروج ، وبخاصّة إذا كان عند بدء الانتقال الى المثلّثة، فإنّه أبلغ في القوّة .}
\EdL{9}{AB02-PDF0018-L09}{}{وبَعْدَه يكون ما قُلنا عند عودة القران إلى البرج الذي ابتدئ فيه من المثلّثة . وبَعْدَه إذا كان}
\EdL{10}{AB02-PDF0018-L10}{}{في البرجين الباقيين من المثلّثة ، وهو أضعف الثلاثة في القوة . وذلك أنّه لا يقع قران وكسوف}
\EdL{11}{AB02-PDF0018-L11}{}{في موضع واحد من فلك البروج عند بدء وانتقال القران إلى المثلّثة إلاّ في الفرط من الزمان .}
\EdLp{12}{AB02-PDF0018-L12}{}{وقد قيل إنّ أوقات التغييرات تكون في كل ثلاثمائة وأربعين سنة ونحوها. وذلك أنّه لا يقع}
\EdL{13}{AB02-PDF0018-L13}{}{قران مع كسوف ، أو على مقابلة في فلك البروج فيما وجدنا من المحنة ، إلاّ في القران الخامس}
\EdL{14}{AB02-PDF0018-L14}{}{من الأوّل فإنّه وقع على قطر الكسوف وذلك من بعد ثمانين سنة بالتقريب ،}
\EdL{15}{AB02-PDF0018-L15}{}{ثم من بعد تحويل القران إلى المثلّثة الأخرى ، وهو القران الثامن عشر من الأوّل وذلك من}
\EdL{16}{AB02-PDF0018-L16}{}{بعد ثلاثمائة وسبعة وثلاثين سنة ، فإنّه أيضًا وقع الكسوف في آخر القران ، ثم الذي}
\EdL{17}{AB02-PDF0018-L17}{}{بعده بقرانين ، وهو قران العشرين من الأوّل ، وذلك بعد ثلاثمائة وسبعة وسبعين سنة ،}
\par\smallskip\noindent\rule{40mm}{0.4pt}\par
\AppL{AB02-PDF0018-A01}{١- الأشياء: أجزاء مطموسة في الأصل ، العصيّ والبوقات : كلمة العصيّ سقطت وقد أضيفت كما في الملحق ، وتقام:}
\AppL{AB02-PDF0018-A02}{ويقام ، قيل: قبل}
\PageHead{20}{32}{0020}
\EdL{1}{AB02-PDF0020-L01}{\RL{\ArBody ٣١و}}{/ فإنّه اتفق فيه الكسوف في برج القران إلاّ أنّه لم يكن فيما حسبنا يقع في موضع القران}
\EdL{2}{AB02-PDF0020-L02}{}{كالأوّل ، فلا يكون زحل والمشتري فيه مقترنين كما كانا في الأصل، بل يكون}
\EdL{3}{AB02-PDF0020-L03}{}{بينهما مقدار حد واحد من حدود الكواكب.}
\EdLh{4}{AB02-PDF0020-L04}{}{ظهور النبيّ صلّى اللّه عليه وسلم}
\EdLp{5}{AB02-PDF0020-L05}{}{كان ظهور نبوّته ، صلّى اللّه عليه ، بعد أربعين سنة من القران الأوّل الذي كان في وقت مولده.}
\EdL{6}{AB02-PDF0020-L06}{}{وقد كان أيضًا من القران الثاني عشر من الأوّل الذي من بعد مائتي وثمان وثلاثين سنة قبالة الكسوف ،}
\EdL{7}{AB02-PDF0020-L07}{}{غير أنّ الشمس كانت في آخر جزء من البرج ، وزحل والمشتري في أكثر من نصف البرج ، والقمر}
\EdL{8}{AB02-PDF0020-L08}{}{يقابل الشمس وهو أضعف ما وجدنا في هذه المدّة التي قسنا بها ما ذكرنا من القرانات .}
\EdL{9}{AB02-PDF0020-L09}{}{وربّما خرج القران في وقت الكسوف عن المثلّثة التي يكون فيها من قبل تمام اثني عشر قرانًا.}
\EdL{10}{AB02-PDF0020-L10}{}{وربّما عاد إليها في القران الثالث عشر الذي هو القران الذي يكون فيه الانتقال عنها. فهذا هو المأخذ}
\EdL{11}{AB02-PDF0020-L11}{}{الذي يُستدلّ به على أمور التغييرات .وبمثله أيضًا يُستدلّ على أمور الدول والملوك إلاّ أنّه}
\EdL{12}{AB02-PDF0020-L12}{}{لا يُضطرّ في أمور الدول أن يكون القران مع الكسوف أ[و] على قطره من فلك البروج .}
\EdL{13}{AB02-PDF0020-L13}{}{وكذلك لا يضطر في أمور الملوك أن تكون الكسوفات في أوقات القرانات ولا في مواضعها}
\EdL{14}{AB02-PDF0020-L14}{}{من فلك البروج . بل يكون النظر في أمور الدول التي تكون في البيوتات من حلول الكواكب ا}
\EdL{15}{AB02-PDF0020-L15}{}{لتي تكون أسباب الحوادث في المواضع الرئيسيّة من الابتداء الأوّل الدالّة على التغييرات}
\EdL{16}{AB02-PDF0020-L16}{}{في أوقات الكسوفات الكائنة عند عودات القرانات إلى مواضعها من فك البروج .}
\EdL{17}{AB02-PDF0020-L17}{}{وذلك شيء يكون في كل ستين سنة بالتقريب مرّة واحدة . وربّما تقدم الابتداء في الدولة}
\par\smallskip\noindent\rule{40mm}{0.4pt}\par
\AppL{AB02-PDF0020-A01}{٣- حد: زحل | ٧ - البرج\textsuperscript{١}: البروج | ٨- المدّة : المرة ، القرانات: القران ، وقد صحّحت على الهامش الأيسر بالعبارة التالية «وقيل}
\AppL{AB02-PDF0020-A02}{القرانات» | ١٠- المأخذ : الأخذ ، كذا في الأصل ، ولكن قرأناها وفقًا للعبارة التالية الواردة على هامش المخطوط الأيمن «وفي نسخة}
\AppL{AB02-PDF0020-A03}{المأخذ» | ١٣- مواضعها: كتبت أوّلاً «مواضع» ثم صحّحت بإضافة الهاء دون شطب دائرة العين | ١٥ - التغييرات: وردت العبارة}
\AppL{AB02-PDF0020-A04}{التالية على الهامش: القران (كذا) العلويين في سنة ١١١٤ هجريّة في الحمل وعاد في سنة ١١٧٥ الهجرية (كذا) في الحمل ويكون}
\AppL{AB02-PDF0020-A05}{بينهما ستين سنة .}
\PageHead{22}{34}{0022}
\EdL{1}{AB02-PDF0022-L01}{\RL{\ArBody ٣١ظ}}{/فكان في القران الذي يقع قبل قران العودة قرب انقضائه . ثم يكون تمام الدولة واستواؤها}
\EdL{2}{AB02-PDF0022-L02}{}{في قران العودة من الحال التي تكون للكواكب بعضها عند بعض في هذه الأوقات المذكورة}
\EdL{3}{AB02-PDF0022-L03}{}{ومشاكلتها للأمر الأوّل ، وإسعادها وإضرارها به على حسب ما سيعود فيما يُستأنف .}
\EdLp{4}{AB02-PDF0022-L04}{}{ومن البروج التي تنتهي إليها السنون من البرج الأوّل الذي يكون الابتداء منه لكلّ برج}
\EdL{5}{AB02-PDF0022-L05}{}{سنة كما يُفعل في المواليد . ويكون النظر في أمور الملوك من سائر الكسوفات التي تحدث قبل أوقات}
\EdL{6}{AB02-PDF0022-L06}{}{ملكهم على ما تقدّم من القول وممّا سيظهر فيما يُستقبل من أسباب تلف الملوك وقيامهم . وينظر معها}
\EdL{7}{AB02-PDF0022-L07}{}{في تحاويل السنين التي تكون لمولد صاحب الملّة وبروج الانتهاء من طالعه في تلك السنين}
\EdL{8}{AB02-PDF0022-L08}{}{ومن كسوفات القرانات التي يقع قيام الملوك فيما بينها . ويُقاس بعض الأمور ببعض .}
\EdLp{9}{AB02-PDF0022-L09}{}{وإذ قد أتينا على ما ينبغي أن يُقدم في هذا النوع ، فإنّا نجعل له مثالاً وقياسًا يُقاس إليه .}
\EdL{10}{AB02-PDF0022-L10}{}{ونبتدئ من القران الذي دلّ على ظهور أمر العرب ومولد النبيّ ، صلّى اللّه عليه وسلّم ، ونلقي سواه ممّا قبله}
\EdL{11}{AB02-PDF0022-L11}{}{إذ كان القول فيما مضى من الأمم السالفة والملل.}
\EdLh{12}{AB02-PDF0022-L12}{}{وكان القران الذي دلّ على أمر العرب وظهور النبيّ ، صلّى اللّه عليه وسلّم، ومولده}
\EdL{13}{AB02-PDF0022-L13}{}{عند ابتداء الانتقال من مثلّثة الميزان إلى مثلّثة السرطان . فكان بدء انتقاله إليها في برج العقرب .}
\EdL{14}{AB02-PDF0022-L14}{}{واتفق معه في وقته وفي موضعه من فلك البروج كسوف القمر أعظم ما يكون من الكسوفات}
\EdL{15}{AB02-PDF0022-L15}{}{وأقربه منه إليها في الليلة التي صبيحتها اليوم السّادس والعشرون من نيسان من سنة ثمان مائة}
\EdL{16}{AB02-PDF0022-L16}{}{واثنتين وثمانين من سنيّ ذي القرنين ، وهي ليلة الأحد . واتفق طالع الكسوف بمكة التي بها كان}
\EdL{17}{AB02-PDF0022-L17}{}{مولده وظهوره وبيعته ، صلّى اللّه عليه وسلّم ، ومواضع الكواكب بحساب الممتحن على ما نذكره}
\par\smallskip\noindent\rule{40mm}{0.4pt}\par
\AppL{AB02-PDF0022-A01}{٣ - مشاكلتها: مشاكلها | ٦- معها: مع | ٩ - وقياسًا: وقياس | ١٥ - والعشرون: والعشرين | ١٦- واثنتين : واثنتي}
\PageHead{24}{36}{0024}
\par\medskip\begin{center}\begin{tikzpicture}[x=12mm,y=12mm,line width=0.5pt]
\draw (0,0) rectangle (9,9); \draw (3,0)--(3,9) (6,0)--(6,9) (0,3)--(9,3) (0,6)--(9,6);
\draw (0,9)--(3,6) (9,9)--(6,6) (0,0)--(3,3) (9,0)--(6,3);
\node[align=center,font=\scriptsize,inner sep=0pt] at (2.0,8.0) {\RL{\ArC الجدي}};
\node[align=center,font=\scriptsize,inner sep=0pt] at (1.0,7.0) {\RL{\ArC الدلو}};
\node[align=center,font=\scriptsize,inner sep=0pt] at (4.5,7.5) {\RL{\ArC القوس و}\\\RL{\ArC حد المشتري}};
\node[align=center,font=\scriptsize,inner sep=0pt] at (6.55,8.45) {\RL{\ArC العقرب}\\\RL{\ArC القمر و ه}\\\RL{\ArC المشتري ب}\\\RL{\ArC زحل ب ل}};
\node[align=center,font=\scriptsize,inner sep=0pt] at (8.0,7.0) {\RL{\ArC الميزان}};
\node[align=center,font=\scriptsize,inner sep=0pt] at (7.5,4.5) {\RL{\ArC السنبلة يد}};
\node[align=center,font=\scriptsize,inner sep=0pt] at (1.5,4.5) {\RL{\ArC الحوت يج}};
\node[align=center,font=\scriptsize,inner sep=0pt] at (8.0,2.0) {\RL{\ArC الأسد}};
\node[align=center,font=\scriptsize,inner sep=0pt] at (6.85,0.85) {\RL{\ArC المرّيخ يد}\\\RL{\ArC السرطان}};
\node[align=center,font=\scriptsize,inner sep=0pt] at (4.5,1.5) {\RL{\ArC الجوزاء و}};
\node[align=center,font=\scriptsize,inner sep=0pt] at (0.85,2.15) {\RL{\ArC الحمل}\\\RL{\ArC عطارد يد ك}};
\node[align=center,font=\scriptsize,inner sep=0pt] at (2.45,0.55) {\RL{\ArC الشمس و ه}\\\RL{\ArC الرأس و ل}\\\RL{\ArC الثور و}\\\RL{\ArC الزهرة ح ي}};
\end{tikzpicture}\end{center}\medskip
\EdL{1}{AB02-PDF0024-L01}{\RL{\ArBody ٣٢و}}{/في هذه الصورة :}
\EdL{2}{AB02-PDF0024-L02}{}{الطالع القوس – بيت المشتري وحدّه ومثلّثته –}
\EdL{3}{AB02-PDF0024-L03}{}{ووسط [ السماء ] السنبلة ، بيت عطارد وشرفه وحدّ الزهرة ،}
\EdL{4}{AB02-PDF0024-L04}{}{ومثلّثته القمر ثم الزهرة . الشمس في}
\EdL{5}{AB02-PDF0024-L05}{}{الثور، حد الزهرة وبيتها وشرف القمر ومثلّثته ،}
\EdL{6}{AB02-PDF0024-L06}{}{ثم مثلّثته الزهرة. والقمر في العقرب ،}
\EdL{7}{AB02-PDF0024-L07}{}{عرضه في الشمال ، ب ، بيت المرّيخ ومثلّثته وآخر حده . زحل والمشتري في العقرب مقترنين .}
\EdL{8}{AB02-PDF0024-L08}{}{أمّا عرض زحل فإنّه د مه وعرض المشتري ا ن . وهما هابطان في الشمال في العرض في حدّ المرّيخ}
\EdL{9}{AB02-PDF0024-L09}{}{وبيته ومثلّثته . المرّيخ في السرطان ، عرضه ا ، هابط في الشمال في حدّ نفسه ومثلّثته،}
\EdL{10}{AB02-PDF0024-L10}{}{وهو برج هبوطه مع الكوكب الذي على رأس التوأمين ، ومزاجه مزاج المرّيخ . والزهرة}
\EdL{11}{AB02-PDF0024-L11}{}{في الثور وعرضها هـ ن في الشمال ، صاعدة في بيتها ومثلّثتها مع القمر وحدّ المشتري . وعطارد في}
\EdL{12}{AB02-PDF0024-L12}{}{الحمل، عرضه ا ن ، في الشمال هابط في حدّ نفسه ، ومثلّثته المشتري . الرأس في الثور . والكوكب}
\EdL{13}{AB02-PDF0024-L13}{}{الذي طلع مع القمر قبل كسوفه من الكواكب الثابتة برجل فرس قنطورس، ومزاجه مزاج}
\EdL{14}{AB02-PDF0024-L14}{}{الزهرة والمشتري . فكان هذا الكسوف والقران على مذهب بطلميوس الذي أوضحناه}
\EdL{15}{AB02-PDF0024-L15}{}{في صدر هذا الكتاب ، أنّه لمّا اجتمع القران والكسوف عند بدء الانتقال في وقت واحد}
\EdL{16}{AB02-PDF0024-L16}{}{وحدّ واحد من العقرب – وهي برج العرب التي تشاكلها مع المرّيخ الذي ينسب إليه من}
\EdL{17}{AB02-PDF0024-L17}{}{بلدانهم الحجاز مع الزهرة ، وخاصّة مدينة يثرب – دلّ على قوة أمر العرب وظهورهم}
\par\smallskip\noindent\rule{40mm}{0.4pt}\par
\AppL{AB02-PDF0024-A01}{٣ - السماء : سقطت في الأصل ، وقد أضيفت هنا لاستقامة المعنى | ١٠ - الكوكب : الكواكب | ١٣ - فرس: تحت السطر}
\AppL{AB02-PDF0024-A02}{عبارة «من الثوابت» | ١٤ - فكان: فكانت | ١٦ - وهي: كذا يستخدم برج على أنّها مؤنّث ، تشاكلها: تشاكلهنَّ .}
\end{document}
